\subsection{All small state sizes in the faithful-hull square question}

\begin{lemma}[Proved: a source-fixed state opposite $R_3$]
\label{CP2:A-hull-square-fixed-state}
Let $N=R_3$, let $M$ be any monoid, and let $S$ be a
nonempty left-$N$, right-$M$ biset.
If some $s\in S$ satisfies $sm=s$ for every $m\in M$,
then all products of degree-two classes of its collage
vanish over every field, including $\mathbb F_2$.
This conclusion has no higher-homology hypothesis.
\end{lemma}
\begin{proof}
In the actual two-sided bridge $E=B(N,S,M)$, select $s$
and use only identity strings on the $N$ side.
The fixed-state equations give a simplicial section
$BM\to E$ of the bridge projection to $BM$.
Its projection to $BN$ is constant.
Since $BR_3$ is contractible, every second class of the
collage has zero restriction to $BN$.
Compatibility of its two corner restrictions in the
Mayer--Vietoris sequence, followed by the section,
forces its restriction to $BM$ to vanish as well.
The two-block relative-product lemma now kills every
product of second classes.  Equivalently, both classes
lift from the relative suspension quotient of the
bridge, whose products in these degrees vanish.
All steps are valid with the specified field coefficients.
\end{proof}

\begin{lemma}[Proved: complete no-fixed-state band acts on at most five states]
\label{CP2:A-three-by-two-small-state-band-forms}
For a fixed labeled $3\times2$ rectangular band, the
numbers of right acts with no common fixed state, up
to state relabeling, are respectively
\[
                    0,\quad1,\quad1,\quad2,\quad5
\]
on $1,2,3,4,5$ states.
If this band is the minimal ideal of a monoid $M$,
having a band-fixed state is equivalent to having
an $M$-fixed state.
\end{lemma}
\begin{proof}
Every minimal-band core orbit is a quotient of the
two column labels.  A one-state orbit is a common
fixed state.  In its absence, every core orbit has
exactly two states, with its column labels distinguished.
On at most five states there is therefore either one
such orbit or two.

With one orbit, each remaining state maps into that
orbit under every band row, and its image under
$e_{ij}$ is the uniquely labeled column-$j$ state.
This gives one act at every state count from two
through five.
With two core orbits, four states are already used.
At four states there is just their disjoint union.
At five states the one remaining state chooses,
independently for each of the three band rows, one
of the two orbits.  Thus it has a tuple in
$\{0,1\}^3$.  Interchanging the two whole core
orbits complements all three bits, giving exactly
four orbits of tuples.  These and the one-core
case give the five forms.
The displayed choices determine every band action,
and every allowed state relabeling acts exactly by
permuting whole core orbits and interchangeable
outside states.  Hence the list is complete.

Finally, if $s$ is fixed by every $k\in K$, choose
one $k\in K$.  For every $m\in M$,
\[
              sm=(sk)m=s(km)=s,
\]
because $km\in K$.  The converse is immediate.
\end{proof}

Consequently a finite check for
Corollary~\ref{CP2:A-three-by-two-eleven-finite-square-cover}
may omit all common-fixed-state inputs by the first
lemma and enumerate the complete lists on two,
three, four and five states from the second lemma.
The one-state case is already covered by the
fixed-state lemma.  This accounts for every possible
size of the actual induced state set.
No proper-retract or third-homology exclusion is
part of this reduction.
